\chapter{Orbitais de fronteira e reatividade}\label{ch:b2:frontier-orbitals}

Um frasco de ciclopentadieno esquecido na prateleira se transforma lentamente
em seu dímero: duas moléculas do mesmo \omterm{def:b2:frontier-orbitals:diels-alder}{dieno} se unem em um composto bicíclico
em ponte, e só um dos estereoisômeros possíveis se forma. Nada nas setas
curvas do volume do 1.º ano prevê qual. Dois orbitais o fazem: o orbital
ocupado mais alto de uma molécula e o orbital vazio mais baixo da outra. Este
capítulo transforma a interação de dois níveis do \cref{ch:b2:diatomic-mos} em
uma ferramenta para a reatividade --- que átomo um nucleófilo ataca, por que
lado, e por que um \omterm{def:b2:frontier-orbitals:diels-alder}{dieno} e um alceno se combinam em um anel em uma única
etapa.

\begin{recall}
\Cref{ch:b2:diatomic-mos}: dois orbitais em interação, a repulsão a quatro
elétrons. \Cref{ch:b2:huckel}: níveis e coeficientes de Hückel dos sistemas
conjugados. O volume do 1.º ano: nucleófilos e eletrófilos, setas curvas,
controle cinético, reações estereoespecíficas, estereosseletivas e
regiosseletivas, adições eletrofílicas aos alcenos, substituição nucleofílica.
\end{recall}

\section{Dois orbitais, de novo}

Quando duas moléculas se aproximam, seus orbitais se sobrepõem e interagem. A
maior parte da interação é fraca, por estarem os orbitais distantes em
energia; basta então uma forma perturbativa do resultado de dois níveis.

\begin{theorem}[Interação fraca de dois níveis]\label{thm:b2:frontier-orbitals:perturbation}
Dois orbitais de energias $\varepsilon_1 < \varepsilon_2$, acoplados por uma \omterm{def:b2:diatomic-mos:integrals}{integral
de ressonância} $\beta$ com $|\beta| \ll \Delta\varepsilon = \varepsilon_2 - \varepsilon_1$
(sobreposição desprezada), são deslocados para
\[
  \varepsilon_1 - \frac{\beta^2}{\Delta\varepsilon}, \qquad \varepsilon_2 + \frac{\beta^2}{\Delta\varepsilon}
\]
em primeira ordem em $\beta^2/\Delta\varepsilon^2$: o nível mais baixo desce e o mais alto
sobe do mesmo valor, tanto menor quanto maior a diferença.
\end{theorem}

\begin{proof}
Pelo \cref{thm:b2:diatomic-mos:heteronuclear}, os níveis exatos são $\bar\varepsilon \mp
\sqrt{\Delta\varepsilon^2/4 + \beta^2}$. Com $u = 4\beta^2/\Delta\varepsilon^2 \ll 1$, $\sqrt{\Delta\varepsilon^2/4
+ \beta^2} = \frac{\Delta\varepsilon}{2}\sqrt{1 + u} \approx \frac{\Delta\varepsilon}{2}(1 + u/2) =
\frac{\Delta\varepsilon}{2} + \frac{\beta^2}{\Delta\varepsilon}$, e $\bar\varepsilon \mp \Delta\varepsilon/2 = \varepsilon_1,
\varepsilon_2$.
\end{proof}

\begin{proposition}[Dois elétrons atraem, quatro repelem]\label{prop:b2:frontier-orbitals:four-electron}
Se os dois orbitais em interação contêm dois elétrons ao todo, o sistema é
estabilizado de cerca de $2\beta^2/\Delta\varepsilon$. Se contêm quatro, o sistema é desestabilizado.
\end{proposition}

\begin{proof}
Dois elétrons vão para o nível abaixado: ganho de $2\beta^2/\Delta\varepsilon$. Com quatro, o nível
superior também é preenchido; na ordem do \cref{thm:b2:frontier-orbitals:perturbation}, os
dois deslocamentos se compensam e, mantida a sobreposição, o nível superior
sobe mais do que o inferior desce (\cref{prop:b2:diatomic-mos:asymmetry}): uma
repulsão líquida.
\end{proof}

\begin{omfigure}
\begin{tikzpicture}[font=\small]
  \begin{scope}
    \pic (a) at (0,0) {omlevel={ud}{0.7}};
    \pic (b) at (3,1.4) {omlevel={empty}{0.7}};
    \pic (lo) at (1.5,-0.5) {omlevel={ud}{0.7}};
    \pic (hi) at (1.5,1.9) {omlevel={empty}{0.7}};
    \draw[omcorr, densely dashed] (a-r) -- (lo-l) (a-r) -- (hi-l) (b-l) -- (lo-r) (b-l) -- (hi-r);
    \node[below=6pt] at (1.5,-0.6) {dois elétrons: estabilizados};
    \node[left] at (a-l) {preenchido};
    \node[right] at (b-r) {vazio};
  \end{scope}
  \begin{scope}[xshift=7cm]
    \pic (a) at (0,0) {omlevel={ud}{0.7}};
    \pic (b) at (3,0.6) {omlevel={ud}{0.7}};
    \pic (lo) at (1.5,-0.5) {omlevel={ud}{0.7}};
    \pic (hi) at (1.5,1.25) {omlevel={ud}{0.7}};
    \draw[omcorr, densely dashed] (a-r) -- (lo-l) (a-r) -- (hi-l) (b-l) -- (lo-r) (b-l) -- (hi-r);
    \node[below=6pt] at (1.5,-0.6) {quatro elétrons: repulsão};
    \node[left] at (a-l) {preenchido};
    \node[right] at (b-r) {preenchido};
  \end{scope}
\end{tikzpicture}

{\small Interação de dois orbitais de moléculas vizinhas. Um orbital
preenchido diante de um vazio dá uma estabilização líquida; dois orbitais
preenchidos se repelem, pois o nível superior sobe mais do que o inferior
desce. Essa repulsão é a origem do impedimento estérico entre moléculas.}
\end{omfigure}

\section{Orbitais de fronteira}

\begin{definition}[Orbitais de fronteira]\label{def:b2:frontier-orbitals:frontier}
O \emph{HOMO}\index{HOMO} (\omterm{def:b2:diatomic-mos:lcao}{orbital molecular} ocupado de maior energia) e o
\emph{LUMO}\index{LUMO} (\omterm{def:b2:diatomic-mos:lcao}{orbital molecular} vazio de menor energia) de uma
molécula são seus \emph{orbitais de fronteira}\index{orbitais de fronteira}.
\end{definition}

\begin{proposition}[A interação dominante]\label{prop:b2:frontier-orbitals:dominant}
Entre duas moléculas de camada fechada, as interações estabilizantes que mais
importam são as entre o \omterm{def:b2:frontier-orbitals:frontier}{HOMO} de uma e o \omterm{def:b2:frontier-orbitals:frontier}{LUMO} da outra; dos dois pares desse
tipo, domina o de menor diferença de energia. Um nucleófilo reage por seu
\omterm{def:b2:frontier-orbitals:frontier}{HOMO}, um eletrófilo por seu \omterm{def:b2:frontier-orbitals:frontier}{LUMO}.
\end{proposition}

\begin{proof}[Raciocínio]
Os pares preenchido--preenchido se repelem (\cref{prop:b2:frontier-orbitals:four-electron}) e os
pares vazio--vazio não contêm elétrons: só os pares preenchido--vazio
estabilizam, cada um de $2\beta^2/\Delta\varepsilon$. O \omterm{def:b2:frontier-orbitals:frontier}{HOMO} é o nível preenchido mais alto e o
\omterm{def:b2:frontier-orbitals:frontier}{LUMO} o vazio mais baixo; logo, um par HOMO--LUMO tem a menor diferença entre
eles; pelo teorema, a menor diferença dá o maior termo.
\end{proof}

\begin{definition}[Controle de carga e controle orbital]\label{def:b2:frontier-orbitals:control}
Uma reação está sob \emph{controle de carga}\index{controle de carga} quando a
atração entre as cargas dos dois parceiros decide onde eles reagem, e sob
\emph{controle orbital}\index{controle orbital} quando é a sobreposição dos \omterm{def:b2:frontier-orbitals:frontier}{orbitais
de fronteira} que decide, ocorrendo o ataque onde seus coeficientes são maiores.
\end{definition}

\begin{definition}[Espécies duras e moles]\label{def:b2:frontier-orbitals:hard-soft}
Um \emph{nucleófilo duro}\index{nucleofilo duro@nucleófilo duro} ou um \emph{eletrófilo duro}
\index{eletrofilo duro@eletrófilo duro} é pequeno, tem uma carga concentrada e um
\omterm{def:b2:frontier-orbitals:frontier}{HOMO} (respectivamente, um \omterm{def:b2:frontier-orbitals:frontier}{LUMO}) distante do de seus parceiros: reage sob
\omterm{def:b2:frontier-orbitals:control}{controle de carga}. Um \emph{nucleófilo mole}\index{nucleofilo mole@nucleófilo mole} ou um \emph{eletrófilo mole}
\index{eletrofilo mole@eletrófilo mole} é grande e polarizável, com um \omterm{def:b2:frontier-orbitals:frontier}{HOMO} alto
(respectivamente, um \omterm{def:b2:frontier-orbitals:frontier}{LUMO} baixo): reage sob \omterm{def:b2:frontier-orbitals:control}{controle orbital}. Duro reage de
preferência com duro, mole com mole.
\end{definition}

O íon hidróxido e o íon fluoreto são \omterm{def:b2:frontier-orbitals:hard-soft}{nucleófilos duros}; o iodeto e os
tiolatos, moles; um próton e um carbocátion são \omterm{def:b2:frontier-orbitals:hard-soft}{eletrófilos duros}; o carbono
de um haloalcano, mole.

\begin{definition}[Nucleófilo ambidentado]\label{def:b2:frontier-orbitals:ambident}
Um \emph{nucleófilo ambidentado}\index{nucleofilo ambidentado@nucleófilo ambidentado} tem dois átomos
nucleofílicos ligados por conjugação, cada um dos quais pode atacar um eletrófilo.
\end{definition}

O íon cianeto é atacado no nitrogênio, onde a carga negativa é maior, pelos
\omterm{def:b2:frontier-orbitals:hard-soft}{eletrófilos duros}, e no carbono, onde seu \omterm{def:b2:frontier-orbitals:frontier}{HOMO} tem o maior coeficiente, pelos
moles: um haloalcano dá uma nitrila \ce{R-C#N}. Os íons enolato, vistos no
\cref{ch:b2:enolates-aldol}, reagem do mesmo modo no carbono ou no oxigênio.

\begin{method}[Uma análise por orbitais de fronteira]\label{met:b2:frontier-orbitals:fmo}
\begin{enumerate}
  \item Identifique o nucleófilo (\omterm{def:b2:frontier-orbitals:frontier}{HOMO} alto) e o eletrófilo (\omterm{def:b2:frontier-orbitals:frontier}{LUMO} baixo).
  \item Desenhe os dois \omterm{def:b2:frontier-orbitals:frontier}{orbitais de fronteira} com seus coeficientes e sinais.
  \item A reação ocorre onde a sobreposição é maior: átomos de maiores
        coeficientes, lobos de mesmo sinal frente a frente.
  \item A direção de aproximação segue a forma do \omterm{def:b2:frontier-orbitals:frontier}{LUMO} (perpendicular a um
        plano \ce{C=O}, pelo lado de trás de uma ligação \ce{C-X}).
  \item Uma diferença HOMO--LUMO menor significa uma reação mais rápida
        (mantidas as demais condições).
\end{enumerate}
\end{method}

A direção de ataque decorre do \omterm{def:b2:frontier-orbitals:frontier}{LUMO}. O \omterm{def:b2:frontier-orbitals:frontier}{LUMO} de um grupo carbonila é seu
orbital $\pi^*$, maior no carbono e perpendicular ao plano: os nucleófilos se
aproximam do carbono por cima ou por baixo do plano, não ao longo do eixo
\ce{C=O}. O \omterm{def:b2:frontier-orbitals:frontier}{LUMO} de um haloalcano é o orbital $\sigma^*$ da ligação \ce{C-X}, cujo
grande lobo no carbono aponta para o lado oposto a X: o nucleófilo ataca por
trás, e o carbono sofre inversão --- a estereoquímica da substituição
bimolecular do volume do 1.º ano.

\section{A reação de Diels--Alder}

\begin{definition}[Reação concertada]\label{def:b2:frontier-orbitals:concerted}
Uma \emph{reação concertada}\index{reacao concertada@reação concertada} forma e rompe todas as
suas ligações em uma única etapa, por um único estado de transição, sem
intermediário.
\end{definition}

\begin{definition}[Reação de Diels--Alder]\label{def:b2:frontier-orbitals:diels-alder}
A \emph{reação de Diels--Alder}\index{reacao de Diels--Alder@reação de Diels--Alder} é a adição concertada
de um \emph{dieno}\index{dieno} conjugado, em sua conformação s-cis, a um alceno ou
alcino, o \emph{dienófilo}\index{dienofilo@dienófilo}, formando um anel de seis
membros com duas novas ligações $\sigma$ e uma nova ligação $\pi$.
\end{definition}

\begin{omfigure}
\begin{tikzpicture}[font=\small, >={Stealth}, scale=1.25]
  % diene s-cis: C1 (0,1.6) C2 (0.6,2.4) C3 (1.6,2.4) C4 (2.2,1.6); dienophile below
  \draw[thick] (0,1.6) -- (0.6,2.4) -- (1.6,2.4) -- (2.2,1.6);
  \draw[thick] (0.2,1.62) -- (0.65,2.22);
  \draw[thick] (2.0,1.62) -- (1.55,2.22);
  \draw[thick] (0.2,0) -- (2.0,0); \draw[thick] (0.35,0.14) -- (1.85,0.14);
  \node[font=\scriptsize, anchor=east] at (-0.05,1.6) {1};
  \node[font=\scriptsize, anchor=south east] at (0.5,2.42) {2};
  \node[font=\scriptsize, anchor=south] at (1.6,2.42) {3};
  \node[font=\scriptsize, anchor=west] at (2.25,1.6) {4};
  % arrows: dienophile pi -> C1..C(dienophile); C1=C2 pi -> C2-C3; C3=C4 pi -> C4-C
  \draw[->, thick, omProp] (1.1,0.22) .. controls (0.9,0.9) and (0.3,0.9) .. (0.12,1.42);
  \draw[->, thick, omProp] (0.45,1.95) .. controls (0.6,2.75) and (1.0,2.8) .. (1.1,2.48);
  \draw[->, thick, omProp] (1.78,1.92) .. controls (2.5,1.6) and (2.4,0.6) .. (2.05,0.2);
  \node at (3.3,1.2) {$\longrightarrow$};
  \begin{scope}[xshift=4.2cm]
    \draw[thick] (0,1.6) -- (0.6,2.4) -- (1.6,2.4) -- (2.2,1.6) -- (2.0,0) -- (0.2,0) -- cycle;
    \draw[thick] (0.72,2.26) -- (1.48,2.26);
  \end{scope}
\end{tikzpicture}

{\small A \omterm{def:b2:frontier-orbitals:diels-alder}{reação de Diels--Alder} do butadieno com o eteno, escrita com setas
curvas: três pares de elétrons se deslocam ao mesmo tempo, em uma etapa
cíclica e concertada; duas ligações $\sigma$ se formam nas extremidades do
\omterm{def:b2:frontier-orbitals:diels-alder}{dieno}, e a ligação $\pi$ passa para seu centro.}
\end{omfigure}

A interação dominante é entre o \omterm{def:b2:frontier-orbitals:frontier}{HOMO} do \omterm{def:b2:frontier-orbitals:diels-alder}{dieno}, $\psi_2$, e o \omterm{def:b2:frontier-orbitals:frontier}{LUMO} do
\omterm{def:b2:frontier-orbitals:diels-alder}{dienófilo}, $\pi^*$. Seus lobos nas extremidades do \omterm{def:b2:frontier-orbitals:diels-alder}{dieno} e nos dois carbonos do
\omterm{def:b2:frontier-orbitals:diels-alder}{dienófilo} têm sinais compatíveis quando o \omterm{def:b2:frontier-orbitals:diels-alder}{dienófilo} se aproxima pela face do
\omterm{def:b2:frontier-orbitals:diels-alder}{dieno}: as duas novas ligações podem se formar ao mesmo tempo, na mesma face de
cada parceiro.

\begin{omfigure}
\begin{tikzpicture}[font=\small]
  \foreach \x/\c in {0/0.602, 1/0.372, 2/-0.372, 3/-0.602} {
    \pic at (\x,2) {omp={1.5*\c}{90}};
  }
  \draw[black!45] (-0.2,2) -- (3.2,2);
  \node[anchor=west] at (3.6,2) {HOMO do dieno ($\psi_2$)};
  \pic at (0,0) {omp={-0.85}{90}}; \pic at (3,0) {omp={0.85}{90}};
  \draw[black!45] (-0.2,0) -- (3.2,0);
  \node[anchor=west] at (3.6,0) {LUMO do dienófilo ($\pi^*$)};
  \draw[omProp, thick, densely dashed] (0,1.45) -- (0,0.6);
  \draw[omProp, thick, densely dashed] (3,1.45) -- (3,0.6);
  \node[font=\scriptsize, omProp, anchor=east] at (-0.15,1.0) {em fase};
  \node[font=\scriptsize, omProp, anchor=west] at (3.15,1.0) {em fase};
\end{tikzpicture}

{\small \omterm{def:b2:frontier-orbitals:frontier}{Orbitais de fronteira} do butadieno e do eteno, com os lobos em
proporção aos coeficientes de Hückel. Frente a frente, os lobos das
extremidades do \omterm{def:b2:frontier-orbitals:frontier}{HOMO} do \omterm{def:b2:frontier-orbitals:diels-alder}{dieno} e os lobos do \omterm{def:b2:frontier-orbitals:frontier}{LUMO} do \omterm{def:b2:frontier-orbitals:diels-alder}{dienófilo} têm os mesmos
sinais nas duas extremidades: as duas ligações $\sigma$ podem se formar juntas,
na mesma face de cada parceiro.}
\end{omfigure}

\begin{proposition}[Estereoespecificidade]\label{prop:b2:frontier-orbitals:da-stereospecific}
A \omterm{def:b2:frontier-orbitals:diels-alder}{reação de Diels--Alder} é estereoespecífica: substituintes cis no \omterm{def:b2:frontier-orbitals:diels-alder}{dienófilo}
ficam cis no produto, substituintes trans continuam trans.
\end{proposition}

\begin{proof}
A reação é concertada e as duas novas ligações se formam na mesma face do
\omterm{def:b2:frontier-orbitals:diels-alder}{dienófilo}: seus dois carbonos nunca giram um em relação ao outro, e as
posições relativas de seus substituintes são levadas para o anel.
\end{proof}

\begin{proposition}[Regiosseletividade]\label{prop:b2:frontier-orbitals:da-regio}
Com parceiros assimétricos, o produto majoritário une o átomo do \omterm{def:b2:frontier-orbitals:diels-alder}{dieno} com o
maior coeficiente em seu \omterm{def:b2:frontier-orbitals:frontier}{HOMO} ao átomo do \omterm{def:b2:frontier-orbitals:diels-alder}{dienófilo} com o maior coeficiente em
seu \omterm{def:b2:frontier-orbitals:frontier}{LUMO}.
\end{proposition}

\begin{proof}[Raciocínio]
No estado de transição, as duas novas ligações ainda não estão igualmente
formadas; a estabilização $2\beta^2/\Delta\varepsilon$ é maior quando os maiores coeficientes se
sobrepõem ($\beta$ cresce com o produto dos coeficientes dos átomos que se
sobrepõem); assim, essa orientação é alcançada com energia menor.
\end{proof}

\begin{omfigure}
\begin{tikzpicture}
\begin{axis}[ybar, width=0.47\linewidth, height=4.6cm, ymin=-0.8, ymax=0.8,
  xtick={0,1,2,3,4}, xticklabels={O,C1,C2,C3,C4}, axis x line=bottom, axis y line=left,
  extra y ticks={0}, extra y tick style={grid=major, grid style={black!50}}, enlarge x limits=0.15,
  tick label style={font=\scriptsize}, ylabel={coeficiente}, label style={font=\small},
  title={\small HOMO do 1-metoxibutadieno}, bar width=8pt]
\addplot[fill=omDef!50, draw=omDef] table[x=atom, y=c] {figdata/out/bachelor-2-huckel-c-diene.dat};
\end{axis}
\end{tikzpicture}\hfill
\begin{tikzpicture}
\begin{axis}[ybar, width=0.47\linewidth, height=4.6cm, ymin=-0.8, ymax=0.8,
  xtick={1,2,3,4}, xticklabels={C$\beta$,C$\alpha$,C(O),O}, axis x line=bottom, axis y line=left,
  extra y ticks={0}, extra y tick style={grid=major, grid style={black!50}}, enlarge x limits=0.15,
  tick label style={font=\scriptsize}, ylabel={coeficiente}, label style={font=\small},
  title={\small LUMO do propenal}, bar width=8pt]
\addplot[fill=omThm!45, draw=omThm] table[x=atom, y=c] {figdata/out/bachelor-2-huckel-c-dienophile.dat};
\end{axis}
\end{tikzpicture}

{\small Modelo de Hückel com parâmetros de heteroátomo ($\alpha_{\ce{O}} = \alpha + 2\beta$,
$\beta_{\ce{CO}} = 0.8\beta$ para o oxigênio do éter; $\alpha_{\ce{O}} = \alpha + \beta$, $\beta_{\ce{CO}} =
\beta$ para o oxigênio da carbonila; um modelo). O \omterm{def:b2:frontier-orbitals:frontier}{HOMO} do \omterm{def:b2:frontier-orbitals:diels-alder}{dieno} é maior em C4, a
extremidade afastada do grupo metoxila (0.60 contra 0.51); o \omterm{def:b2:frontier-orbitals:frontier}{LUMO} do \omterm{def:b2:frontier-orbitals:diels-alder}{dienófilo} é
maior no carbono $\beta$ (0.66). Eles se ligam entre si: os grupos metoxila e
aldeído acabam em carbonos vizinhos do anel.}
\end{omfigure}

O grupo doador eleva o \omterm{def:b2:frontier-orbitals:frontier}{HOMO} do \omterm{def:b2:frontier-orbitals:diels-alder}{dieno} ($m = 0.48$ em vez de 0.62) e o aceptor
abaixa o \omterm{def:b2:frontier-orbitals:frontier}{LUMO} do \omterm{def:b2:frontier-orbitals:diels-alder}{dienófilo} ($m = -0.35$ em vez de $-1$): a diferença diminui, e
esses pares reagem muito mais depressa que o butadieno com o eteno. Um \omterm{def:b2:frontier-orbitals:diels-alder}{dieno}
rico em elétrons e um \omterm{def:b2:frontier-orbitals:diels-alder}{dienófilo} pobre em elétrons formam a combinação típica.

\begin{definition}[Adutos endo e exo]\label{def:b2:frontier-orbitals:endo}
Quando um \omterm{def:b2:frontier-orbitals:diels-alder}{dieno} cíclico se adiciona a um \omterm{def:b2:frontier-orbitals:diels-alder}{dienófilo} que leva um substituinte
insaturado, o \emph{aduto endo}\index{aduto endo} tem esse substituinte voltado
para a ligação dupla recém-formada (sob o \omterm{def:b2:frontier-orbitals:diels-alder}{dieno} no estado de transição), e o
\emph{aduto exo}\index{aduto exo}, voltado para o lado oposto. A \emph{regra
endo}\index{regra endo} afirma que o aduto endo se forma mais depressa.
\end{definition}

\begin{proposition}[A regra endo]\label{prop:b2:frontier-orbitals:endo}
Sob controle cinético, o \omterm{def:b2:frontier-orbitals:endo}{aduto endo} é o produto majoritário, embora o \omterm{def:b2:frontier-orbitals:endo}{aduto exo}
seja frequentemente o mais estável.
\end{proposition}

\admitted

A explicação usual é uma sobreposição secundária, na aproximação endo, entre os
lobos do sistema $\pi$ do substituinte e os dos átomos centrais do \omterm{def:b2:frontier-orbitals:frontier}{HOMO} do
\omterm{def:b2:frontier-orbitals:diels-alder}{dieno}; ela abaixa o estado de transição endo sem formar ligação. A regra é
empírica e tem exceções; o volume do 3.º ano trata essas reações pela simetria
de seus orbitais.

\begin{omfigure}
\begin{tikzpicture}[font=\small]
  \foreach \xs/\mode in {0/endo, 5/exo} {
    \begin{scope}[xshift=\xs cm]
      % cyclopentadiene seen from the side: C1, C2=C3, C4, CH2 bridge on top
      \draw[thick] (0,2) -- (0.6,2.5) -- (1.8,2.5) -- (2.4,2);
      \draw[thick] (0.75,2.38) -- (1.65,2.38);
      \draw[thick] (0,2) -- (1.2,3.2) -- (2.4,2);
      \node[font=\scriptsize, anchor=south] at (1.2,3.2) {\ce{CH2}};
      % dienophile
      \draw[thick] (0.2,0) -- (2.2,0); \draw[thick] (0.35,0.1) -- (2.05,0.1);
      \draw[omProp, densely dashed, thick] (0,1.9) -- (0.2,0.15);
      \draw[omProp, densely dashed, thick] (2.4,1.9) -- (2.2,0.15);
      \node[anchor=north] at (1.2,-1.3) {\mode};
    \end{scope}
  }
  % endo: carbonyls tucked under the diene
  \draw[thick] (0.2,0) -- (0.75,0.95) (2.2,0) -- (1.65,0.95);
  \node[font=\scriptsize, fill=white, inner sep=1pt] at (0.75,1.1) {C=O};
  \node[font=\scriptsize, fill=white, inner sep=1pt] at (1.65,1.1) {C=O};
  \draw[omThm, densely dotted, thick] (0.75,1.25) -- (0.65,2.35);
  \draw[omThm, densely dotted, thick] (1.65,1.25) -- (1.75,2.35);
  % exo: carbonyls pointing away
  \draw[thick] (5.2,0) -- (4.95,-0.85) (7.2,0) -- (7.45,-0.85);
  \node[font=\scriptsize] at (4.95,-1.0) {C=O};
  \node[font=\scriptsize] at (7.45,-1.0) {C=O};
\end{tikzpicture}

{\small Aproximações endo e exo do anidrido maleico ao ciclopentadieno,
esquemáticas. As ligações em formação estão tracejadas. Na aproximação endo,
os grupos carbonila ficam sob o \omterm{def:b2:frontier-orbitals:diels-alder}{dieno}, e seus orbitais $\pi$ se sobrepõem aos
carbonos centrais dele (pontilhado, sobreposição secundária); na aproximação
exo, eles apontam para o lado oposto.}
\end{omfigure}

\begin{method}[Desenhar um produto de Diels--Alder]\label{met:b2:frontier-orbitals:diels-alder}
\begin{enumerate}
  \item Coloque o \omterm{def:b2:frontier-orbitals:diels-alder}{dieno} em sua conformação s-cis, com o \omterm{def:b2:frontier-orbitals:diels-alder}{dienófilo} diante de
        suas extremidades.
  \item Desenhe as duas novas ligações $\sigma$ das extremidades do \omterm{def:b2:frontier-orbitals:diels-alder}{dieno} aos
        carbonos do \omterm{def:b2:frontier-orbitals:diels-alder}{dienófilo}, e a nova ligação $\pi$ entre C2 e C3 do \omterm{def:b2:frontier-orbitals:diels-alder}{dieno}.
  \item Mantenha os substituintes do \omterm{def:b2:frontier-orbitals:diels-alder}{dienófilo} do mesmo lado (cis continua cis).
  \item Para parceiros assimétricos, una os maiores coeficientes (produtos
        ``orto'' ou ``para'' para \omterm{def:b2:frontier-orbitals:diels-alder}{dienos} substituídos em 1 ou em 2).
  \item Com um \omterm{def:b2:frontier-orbitals:diels-alder}{dieno} cíclico, coloque o substituinte insaturado em endo.
\end{enumerate}
\end{method}

\begin{inthelab}[O craqueamento do diciclopentadieno]
O ciclopentadieno é vendido sob a forma de seu dímero e regenerado logo antes
do uso: o dímero é aquecido até seu ponto de ebulição (cerca de
\qty{170}{\degreeCelsius}) em um balão provido de uma coluna de
fracionamento, onde ele se decompõe de volta no monômero (uma reação de
retro-Diels--Alder), que destila a \qty{41}{\degreeCelsius} e é recolhido em um
frasco coletor resfriado em gelo. O monômero torna a dimerizar em poucas horas
à temperatura ambiente e é usado imediatamente. Os dois compostos são
inflamáveis e nocivos; o trabalho é feito em uma capela.
\end{inthelab}

\begin{safety}{\ghs[5mm]{GHS02}\,\ghs[5mm]{GHS06}\,\ghs[5mm]{GHS08}}
O ciclopentadieno e seu dímero são altamente inflamáveis e tóxicos por inalação;
o anidrido maleico é corrosivo e um sensibilizante respiratório. Eles são
manipulados na capela, com luvas e proteção ocular, longe de chamas.
% ledger: ghs:cyclopentadiene, ghs:dicyclopentadiene, ghs:maleic-anhydride, bp:cyclopentadiene, bp:dicyclopentadiene
\end{safety}

\section{Exercícios}

\begin{exercise}[$\star$]\label{exo:b2:frontier-orbitals:1}
Dê o \omterm{def:b2:frontier-orbitals:frontier}{HOMO} e o \omterm{def:b2:frontier-orbitals:frontier}{LUMO} (energias em unidades de Hückel) do eteno, do ânion alila e
do butadieno.
\end{exercise}

\begin{exercise}[$\star$]\label{exo:b2:frontier-orbitals:2}
Classifique como duro ou mole: \ce{OH-}, \ce{I-}, \ce{H+}, um tiolato \ce{RS-}, o
carbono do iodometano, \ce{F-}.
\end{exercise}

\begin{exercise}[$\star$]\label{exo:b2:frontier-orbitals:3}
Quais destes \omterm{def:b2:frontier-orbitals:diels-alder}{dienos} podem participar de uma \omterm{def:b2:frontier-orbitals:diels-alder}{reação de Diels--Alder}:
buta-1,3-dieno, ciclopentadieno, penta-1,4-dieno, (2\textit{Z},4\textit{Z})-hexa-2,4-dieno em sua
forma s-cis?
\end{exercise}

\begin{exercise}[$\star$]\label{exo:b2:frontier-orbitals:4}
Desenhe o produto do butadieno com a propenonitrila \ce{CH2=CH-C#N}.
\end{exercise}

\begin{exercise}[$\star\star$]\label{exo:b2:frontier-orbitals:5}
Dois orbitais em $-10$ e \qty{-2}{eV} interagem com $\beta = \qty{-1.0}{eV}$ (dados do exercício).
Calcule a estabilização perturbativa de dois elétrons e depois a exata.
\end{exercise}

\begin{exercise}[$\star\star$]\label{exo:b2:frontier-orbitals:6}
O íon cianeto reage com o iodometano dando etanonitrila \ce{CH3-C#N}, mas, com um
carbocátion em um solvente fortemente ionizante, dá um pouco de isonitrila
\ce{R-N#C}. Explique.
\end{exercise}

\begin{exercise}[$\star\star$]\label{exo:b2:frontier-orbitals:7}
Usando os coeficientes do capítulo, preveja o produto majoritário do
1-metoxibutadieno com o propenal.
\end{exercise}

\begin{exercise}[$\star\star$]\label{exo:b2:frontier-orbitals:8}
Desenhe o \omterm{def:b2:frontier-orbitals:endo}{aduto endo} do ciclopentadieno com o propenoato de metila.
\end{exercise}

\begin{exercise}[$\star\star$]\label{exo:b2:frontier-orbitals:9}
Por que o anidrido maleico reage com o ciclopentadieno à temperatura ambiente,
enquanto o eteno exige aquecimento sob pressão?
\end{exercise}

\begin{exercise}[$\star\star\star$]\label{exo:b2:frontier-orbitals:10}
O ciclopentadieno reage com o maleato de dimetila (diéster cis) e com o
fumarato de dimetila (diéster trans). Desenhe os dois produtos e diga quantos
estereoisômeros cada um dá.
\end{exercise}

\begin{exercise}[$\star\star\star$]\label{exo:b2:frontier-orbitals:11}
A \omterm{def:b2:frontier-orbitals:diels-alder}{reação de Diels--Alder} tem $\Delta_r H^\circ < 0$ e $\Delta_r S^\circ < 0$. Explique os
sinais e por que a reação inversa se torna favorável a alta temperatura.
\end{exercise}

\begin{exercise}[$\star\star\star$]\label{exo:b2:frontier-orbitals:12}
Na hidrazina \ce{H2N-NH2}, cada nitrogênio tem um par não ligante. Usando a
interação a quatro elétrons, explique por que a molécula evita a conformação
em que os dois pares não ligantes são paralelos.
\end{exercise}

\section{Problema: O craqueamento do ciclopentadieno}

\begin{problem}[{Problema de fim de semana --- a dimerização do ciclopentadieno
como reação de Diels--Alder, a termodinâmica do craqueamento, os orbitais de
fronteira do ciclopentadieno e do anidrido maleico, e a estereoquímica de seu
aduto}]
\label{pb:b2:frontier-orbitals:1}
Dados: pontos de ebulição, ciclopentadieno \qty{41}{\degreeCelsius},
diciclopentadieno cerca de \qty{170}{\degreeCelsius}. Energias do modelo de
Hückel (um modelo): ciclopentadieno, tratado como uma unidade butadieno, \omterm{def:b2:frontier-orbitals:frontier}{HOMO}
$\alpha + 0.62\beta$ e \omterm{def:b2:frontier-orbitals:frontier}{LUMO} $\alpha - 0.62\beta$; anidrido maleico, \omterm{def:b2:frontier-orbitals:frontier}{LUMO}
cerca de $\alpha - 0.35\beta$.
% ledger: bp:cyclopentadiene, bp:dicyclopentadiene

\medskip
\noindent\textbf{Parte I --- O dímero.}
\begin{enumerate}
  \item Desenhe o ciclopentadieno e mostre que sua unidade \omterm{def:b2:frontier-orbitals:diels-alder}{dieno} é mantida s-cis.
  \item Na dimerização, uma molécula é o \omterm{def:b2:frontier-orbitals:diels-alder}{dieno}. O que é a outra?
  \item Desenhe o dímero e numere as novas ligações.
  \item Que aduto, endo ou exo, se forma à temperatura ambiente, e por quê?
  \item A dimerização é estereoespecífica? Explique pelo mecanismo.
  \item Por que a dimerização não precisa de catalisador?
\end{enumerate}

\medskip
\noindent\textbf{Parte II --- A termodinâmica do craqueamento.}
\begin{enumerate}[resume]
  \item Dê o sinal do $\Delta_r H^\circ$ da dimerização a partir das ligações
        formadas e rompidas.
  \item Dê o sinal do $\Delta_r S^\circ$.
  \item Escreva $\Delta_r G^\circ(T)$ e mostre que ele muda de sinal em uma temperatura $T_i$.
  \item Por que o dímero é estável à temperatura ambiente, e o monômero é favorecido a quente?
  \item Na montagem de craqueamento, por que destilar o monômero à medida que
        ele se forma desloca a reação?
  \item Por que o monômero deve ser mantido frio e usado depressa?
\end{enumerate}

\medskip
\noindent\textbf{Parte III --- \omterm{def:b2:frontier-orbitals:frontier}{Orbitais de fronteira}.}
\begin{enumerate}[resume]
  \item Para o ciclopentadieno reagindo consigo mesmo, calcule a diferença
        HOMO--LUMO em unidades de $|\beta|$.
  \item Para o ciclopentadieno com o anidrido maleico, calcule a diferença entre
        o \omterm{def:b2:frontier-orbitals:frontier}{HOMO} do \omterm{def:b2:frontier-orbitals:diels-alder}{dieno} e o \omterm{def:b2:frontier-orbitals:frontier}{LUMO} do anidrido.
  \item Que par reage mais depressa? Explique.
  \item Por que os dois grupos carbonila do anidrido maleico abaixam seu \omterm{def:b2:frontier-orbitals:frontier}{LUMO}?
  \item Qual orbital do anidrido maleico dá a sobreposição secundária da
        aproximação endo?
\end{enumerate}

\medskip
\noindent\textbf{Parte IV --- O aduto com o anidrido maleico.}
\begin{enumerate}[resume]
  \item Desenhe o \omterm{def:b2:frontier-orbitals:endo}{aduto endo}.
  \item Marque seus estereocentros.
  \item O aduto é quiral? (Procure um plano de simetria.)
  \item Por que os dois hidrogênios do antigo \omterm{def:b2:frontier-orbitals:diels-alder}{dienófilo} ficam do mesmo lado do anel?
  \item Como seria o \omterm{def:b2:frontier-orbitals:endo}{aduto exo}?
  \item Os \omterm{def:b2:frontier-orbitals:endo}{adutos endo} e exo podem se interconverter à temperatura ambiente? Por quê?
  \item Como o anidrido reagiria com a água, e que produto se formaria?
  \item Dê o número de estereocentros do \omterm{def:b2:frontier-orbitals:endo}{aduto endo} do ciclopentadieno com o
        anidrido maleico.
\end{enumerate}
\end{problem}
